% ============================================================
% ICBHI 2017 Lung Sound Classification Paper
% Auto-generated by Paper Writer Agent
% ============================================================

\documentclass[conference]{IEEEtran}

\usepackage{cite}
\usepackage{amsmath,amssymb,amsfonts}
\usepackage{algorithmic}
\usepackage{graphicx}
\usepackage{textcomp}
\usepackage{xcolor}
\usepackage{booktabs}
\usepackage{multirow}
\usepackage{subcaption}
\usepackage{hyperref}

\definecolor{myblue}{RGB}{0,102,204}

\begin{document}

\title{ {{TITLE}} }

\author{
    \IEEEauthorblockN{ {{AUTHOR_LIST}} }
    \IEEEauthorblockA{
        \textit{ {{AFFILIATION}} } \\
        \textit{ {{DEPARTMENT}} } \\
        \{ {{EMAILS}} \}
    }
}

\maketitle

% ============================================================
% ABSTRACT
% ============================================================
\begin{abstract}
{{ABSTRACT_TEXT}}

\textbf{Results:} Our method achieves a Score of \textbf{ {{SOTA_SCORE}}\% } 
(Se: {{SOTA_SE}}\%, Sp: {{SOTA_SP}}\%) on the official ICBHI 2017 60/40 split, 
establishing a new state-of-the-art for per respiratory cycle lung sound classification.

\textbf{Code:} \url{https://github.com/{{REPO}}}
\end{abstract}

% ============================================================
% INTRODUCTION
% ============================================================
\section{Introduction}

Lung auscultation is one of the most common diagnostic procedures in respiratory medicine. 
Abnormal lung sounds such as wheezes and crackles are key indicators of various pulmonary 
conditions including asthma, COPD, pneumonia, and bronchitis~\cite{pasterkamp1997respiratory}. 
However, manual auscultation suffers from high inter-rater variability and requires 
significant clinical expertise.

The ICBHI 2017 Challenge~\cite{rocha2019icbhi} provided the first large-scale benchmark 
dataset for automated lung sound classification, containing 6,898 respiratory cycles 
annotated across 126 patients. Despite significant progress, current state-of-the-art 
methods still achieve moderate performance. The best verified Score on the
official 60/40 split should be confirmed through systematic literature review.

{{CONTRIBUTION_PARAGRAPHS}}

Our contributions are:
\begin{itemize}
    \item {{CONTRIBUTION_1}}
    \item {{CONTRIBUTION_2}}
    \item {{CONTRIBUTION_3}}
    \item {{CONTRIBUTION_4}}
\end{itemize}

% ============================================================
% RELATED WORK
% ============================================================
\section{Related Work}

\subsection{Traditional Lung Sound Classification}
Early approaches used hand-crafted features such as Mel-frequency cepstral coefficients 
(MFCCs), wavelet transforms, and spectral features combined with SVM or Random Forest 
classifiers~\cite{pramono2017automatic}.

\subsection{Deep Learning Approaches}
CNN-based methods~\cite{aykanat2017classification} improved significantly over traditional 
methods. Gairola et al.~\cite{gairola2021respirenet} proposed RespireNet using Bi-LSTM 
with attention, achieving 65.93\% Score. Multi-scale architectures and hybrid CNN-RNN 
models~\cite{perna2022deep} further pushed performance to 66.70\%.

\subsection{Pretrained Models and Transformers}
The adoption of Audio Spectrogram Transformer (AST)~\cite{gong2021ast} marked a turning 
point. Recent directions include medical audio pretraining, contrastive learning,
and self-supervised methods — though exact ICBHI 2017 scores for these approaches
must be verified from published sources before citing. The Research Agent should
conduct a systematic literature review to identify the current true SOTA.

{{ADDITIONAL_RELATED_WORK}}

% ============================================================
% PROPOSED METHOD
% ============================================================
\section{Proposed Method}

\subsection{Overview}
{{METHOD_OVERVIEW}}

\subsection{Backbone Architecture}
{{BACKBONE_DESCRIPTION}}

\subsection{ {{INNOVATION_NAME}} }
{{INNOVATION_DESCRIPTION}}

\subsection{Training Strategy}
{{TRAINING_STRATEGY}}

% ============================================================
% EXPERIMENTAL SETUP
% ============================================================
\section{Experimental Setup}

\subsection{Dataset}
We use the ICBHI 2017 dataset~\cite{rocha2019icbhi}, which contains 920 audio recordings 
from 126 patients, totaling approximately 5.5 hours. Respiratory cycles are annotated with 
four classes: Normal, Wheeze, Crackles, and Both (wheezes + crackles simultaneously).

\subsection{Official Split}
Following the official protocol, we split patients 60/40 for training and testing, ensuring 
no patient appears in both sets. This yields {{TRAIN_CYCLES}} training cycles and 
{{TEST_CYCLES}} test cycles.

\subsection{Preprocessing}
{{PREPROCESSING_DETAILS}}

\subsection{Implementation Details}
{{IMPLEMENTATION_DETAILS}}

% ============================================================
% RESULTS
% ============================================================
\section{Results and Discussion}

\subsection{Main Results}
% AUTO-GENERATED TABLE
\begin{table}[h]
\centering
\caption{Comparison with state-of-the-art on ICBHI 2017 official 60/40 split.}
\label{tab:sota}
\begin{tabular}{lccc}
\toprule
\textbf{Method} & \textbf{Se (\%)} & \textbf{Sp (\%)} & \textbf{Score (\%)} \\
\midrule
{{SOTA_TABLE_ROWS}}
\midrule
\textbf{Ours} & \textbf{ {{OUR_SE}} } & \textbf{ {{OUR_SP}} } & \textbf{ {{OUR_SCORE}} } \\
\bottomrule
\end{tabular}
\end{table}

\subsection{Per-Class Analysis}
% AUTO-GENERATED TABLE
\begin{table}[h]
\centering
\caption{Per-class performance of our method.}
\label{tab:perclass}
\begin{tabular}{lcccc}
\toprule
\textbf{Class} & \textbf{Se (\%)} & \textbf{Sp (\%)} & \textbf{Score (\%)} & \textbf{F1 (\%)} \\
\midrule
{{PER_CLASS_ROWS}}
\bottomrule
\end{tabular}
\end{table}

\subsection{Confusion Matrix}
\begin{figure}[h]
\centering
\includegraphics[width=\columnwidth]{figures/confusion_matrix.pdf}
\caption{Confusion matrix on the ICBHI 2017 test set.}
\label{fig:cm}
\end{figure}

\subsection{Discussion}
{{DISCUSSION_TEXT}}

% ============================================================
% ABLATION STUDIES
% ============================================================
\section{Ablation Studies}

{{ABLATION_TEXT}}

% ============================================================
% CONCLUSION
% ============================================================
\section{Conclusion}

{{CONCLUSION_TEXT}}

Future work includes {{FUTURE_WORK}}.

% ============================================================
% REFERENCES
% ============================================================
\bibliographystyle{IEEEtran}
\bibliography{references}

\end{document}
